% CLAIM: none of ours.
\section{Related work}
\label{sec:related}

\paragraph{ES at scale.} \citet{salimans2017} distributed ES across many workers by having
them exchange only scalar returns rather than gradients; ARS \citep{mania2018} and deep
neuroevolution \citep{such2017,conti2017,chrabaszcz2018} use the same design.
\citet{qiu2025} and EGGROLL \citep{eggroll2025} extend it to language models with the two
perturbation designs compared here: full-rank noise regenerated from seeds, and low-rank
noise. Neither compares update placements, and neither implements both designs.

\paragraph{Tuning language models.} MeZO \citep{malladi2023} also estimates gradients from
forward passes, but to train within the memory needed for inference rather than to evaluate a
population in parallel; its constraint is memory, not the communication and repeated
computation measured here. FedKSeed \citep{qin2024fedkseed} communicates only scalar
coefficients in federated full-parameter tuning, and each client reconstructs the update from
shared random seeds, as in replication. Ferret \citep{shu2025ferret} uses shared randomness to
communicate projected first-order updates. In both, the cost of reconstruction matters,
although their optimization and communication schedules differ from a synchronous ES
population. ESSA \citep{korotyshova2025essa} applies ES to align models of up to 72B
parameters by searching a small subspace of LoRA singular values; it does not study
full-parameter update placement at those sizes.

\paragraph{Other ES variants.} Natural ES, CMA-ES and information-geometric optimization adapt
richer search distributions \citep{wierstra2014,hansen2001,hansen2016,ollivier2017}; here the
update rule is held fixed and only the perturbation design varies.

\paragraph{Distributed systems.} Spending extra computation or memory to reduce communication
is a classical systems tradeoff; replicating data in matrix multiplication is one example
\citep{solomonik2011}. This paper measures how the perturbation design changes that tradeoff
for reconstructing the ES update. evosax \citep{lange2022evosax} provides ES algorithms in JAX.
\citet{sheng2020tpu} run distributed ES on TPUs for meta-learning, without comparing
placements. ZeRO \citep{rajbhandari2020zero} removes repeated optimizer
work and memory in data-parallel training by dividing the optimizer step among devices, much
as splitting divides reconstruction; its reduce-scatter and all-gather move the same bytes as
the all-reduce they replace. BOOST \citep{wang2026boost} places the tensor-parallel collective
of a factorized layer at its narrow activation bottleneck, of width $r$ rather than the hidden
size. Splitting, by contrast, communicates full partial updates; communicating their factors
would be another placement, outside the scope of this paper. PowerSGD \citep{vogels2019}
reduces communication by approximating gradients at low rank, whereas the low-rank
perturbations here still produce a full update to communicate: the two use low rank for
different ends.
